\section{Training Details}
\label{app:training}

\subsection{LoRA Configuration}

LoRA modules are applied to the following layers of InternVL2.5-8B:
attention projections ($W_{QKV}$, $W_O$), MLP layers (gate, down, up),
and the vision--language projector ($\texttt{mlp1.1}$, $\texttt{mlp1.3}$).
Rank $r = 8$, scaling factor $\alpha = 16$, dropout $0.05$, bias disabled.

\subsection{Hyperparameters}

\begin{table}[h]
\centering
\small
\caption{Training hyperparameters by phase.}
\begin{tabular}{lccc}
\toprule
 & Phase 1 (SFT) & Phase 3 (CoT-SFT) & Phase 5 (ADPO) \\
\midrule
Epochs & 2 & 2 & 1 \\
Effective batch size & 8 & 8 & 4 \\
Learning rate & $2 \times 10^{-4}$ & $2 \times 10^{-4}$ & $5 \times 10^{-6}$ \\
Warmup steps & 200 & 100 & 3\% of total \\
Max grad norm & 1.0 & 1.0 & 1.0 \\
Weight decay & 0.0 & 0.0 & 0.0 \\
Precision & bfloat16 & bfloat16 & bfloat16 \\
Optimizer & AdamW & AdamW & AdamW \\
\midrule
DPO $\beta$ & --- & --- & 0.1 \\
DPO $\alpha$ (anchor) & --- & --- & 0.5 \\
\bottomrule
\end{tabular}
\end{table}

\subsection{Input Format}

Each training example is formatted as a multi-turn conversation.
The user turn presents 8 video frames followed by the question text and
lettered options (A--H).
The assistant turn contains the target response:
\begin{itemize}
  \item \textbf{Phase 1 (direct):} \texttt{B) person in green shirt performed punch on person in light pants}
  \item \textbf{Phase 3 (CoT):} \texttt{Reasoning: \{chain\} $\backslash$n Answer: B) person in green shirt...}
\end{itemize}

Option order is shuffled independently per example at training time to
prevent positional bias. At evaluation time, option order is fixed for
cross-checkpoint comparability.

\subsection{CoT Generation Details}

The teacher model (base InternVL2.5-8B, no adapter) receives the 8 video
frames and a structured prompt requesting a 5-step reasoning chain:
(1)~identify all visible people and their apparent roles,
(2)~describe the physical actions occurring,
(3)~note the environment and bystanders,
(4)~determine which answer option matches the visual evidence,
(5)~commit to a final lettered answer.
Generation uses nucleus sampling with temperature $0.7$ and $p = 0.9$.

In Stage~2 (annotation-corrected refinement), the Stage~1 chain, the
ground-truth annotation fields, and the correct answer are fed back to the
teacher, which refines the chain to lead logically to the correct answer
while preserving accurate visual descriptions.

A chain is accepted only if its final 80 characters contain a standalone
letter (A--H) matching the correct answer index.

\subsection{Hardness Taxonomy}

The 10-level distractor hardness taxonomy, ordered from hardest to easiest:

\begin{enumerate}
\item \textbf{Role reversal} --- aggressor and victim swapped; both appear in
  the video.
\item \textbf{Wrong action} --- correct cast, different action verb.
\item \textbf{Wrong victim} --- correct aggressor and action, victim replaced
  with absent person.
\item \textbf{Wrong aggressor} --- correct victim and action, aggressor
  replaced.
\item \textbf{Bystander substitution} --- a bystander from the same video
  fills an active role.
\item \textbf{Wrong location} --- correct cast, different environment.
\item \textbf{Wrong category} --- taxonomically different action or
  environment.
\item \textbf{None claim} --- negation distractor (``no aggressive action'').
\item \textbf{Other in cast} --- mentions a cast member without fitting
  above categories.
\item \textbf{Cross-video} --- people borrowed from an entirely different
  video.
\end{enumerate}

\subsection{Evaluation Protocol}

All three adapter checkpoints are evaluated on the held-out test split
(80\% of videos, 10,581 questions) using greedy decoding
(\texttt{do\_sample=False}, \texttt{max\_new\_tokens=64}).
The model's response is parsed via a letter-extraction regex
(\texttt{$\backslash$b([A-H])$\backslash$b}); if no valid letter is found,
a substring match against option texts is attempted.
The letter-parsed rate (fraction of responses where the regex extracts a
valid letter) serves as a diagnostic of format compliance.
